\documentclass[11pt]{article}
\usepackage[T1]{fontenc}
\usepackage[utf8]{inputenc}
\usepackage{lmodern}
\usepackage[french]{babel}
\usepackage[margin=1in]{geometry}
\usepackage{enumitem}

\begin{document}

\section*{Cas d'utilisation : Inscrire un utilisateur}

\begin{description}[leftmargin=4.2cm, style=nextline]
    \item[Titre de cas d'utilisation] Inscrire un utilisateur

    \item[Acteurs]
    \begin{itemize}[leftmargin=*, nosep]
        \item \textbf{Acteur principal} : Nouvel utilisateur (soumet la demande)
        \item \textbf{Acteur secondaire} : Administrateur ou Administrateur IT
              (traite la demande)
    \end{itemize}

    \item[Description] Tout nouvel utilisateur (Médecin, Administrateur ou
    Administrateur IT) soumet une demande d'inscription. Si les champs saisis
    sont conformes, la demande est enregistrée et transmise à l'Administrateur
    ou à l'Administrateur IT, qui peut l'approuver, la refuser ou la supprimer.
    Ce processus est identique quel que soit le rôle du compte à créer.

    \item[Préconditions]
    \begin{itemize}[leftmargin=*, nosep]
        \item Le système est en cours d'exécution.
        \item L'Administrateur ou l'Administrateur IT est authentifié pour
              traiter les demandes en attente.
    \end{itemize}

    \item[Scénario principal]
    \begin{enumerate}[leftmargin=*, topsep=0pt, itemsep=0pt]
        \item Le nouvel utilisateur accède au formulaire d'inscription.
        \item Le système affiche le formulaire.
        \item L'utilisateur remplit les champs (nom, prénom, identifiant,
              mot de passe, rôle) et soumet le formulaire.
    \end{enumerate}

    \item[Alt : champs conformes]
    \begin{enumerate}[leftmargin=*, topsep=0pt, itemsep=0pt, start=4]
        \item Le système envoie la demande au serveur.
        \item Le serveur enregistre la demande avec le statut «~en attente~».
        \item Le serveur confirme l'enregistrement.
        \item Le système confirme la transmission à l'interface.
        \item L'interface notifie l'utilisateur~:
              «~Demande en attente d'approbation~».
    \end{enumerate}

    L'Administrateur ou l'Administrateur IT consulte les demandes en attente
    et choisit l'une des trois actions suivantes~:

    \textbf{Alt : Admin approuve la demande}
    \begin{itemize}[leftmargin=*, nosep]
        \item Le système crée le compte utilisateur avec le rôle demandé.
        \item Le serveur enregistre le compte dans la base de données.
        \item Le système confirme~: «~Compte créé~».
    \end{itemize}

    \textbf{Alt : Admin refuse la demande}
    \begin{itemize}[leftmargin=*, nosep]
        \item Le système passe la demande au statut «~refusée~».
    \end{itemize}

    \textbf{Alt : Admin supprime la demande}
    \begin{itemize}[leftmargin=*, nosep]
        \item L'Administrateur demande la suppression.
        \item Le serveur supprime définitivement la demande de la base de données.
    \end{itemize}

    \item[Alt : champs non conformes]
    \begin{itemize}[leftmargin=*, nosep]
        \item Si un champ est vide, si l'identifiant contient des caractères
              non autorisés, ou si le mot de passe ne respecte pas les règles
              de complexité~: le système met en évidence les erreurs
              et bloque la soumission.
    \end{itemize}

    \item[Postconditions]
    \begin{itemize}[leftmargin=*, nosep]
        \item \textbf{Approbation}~: le compte est actif et l'utilisateur peut
              s'authentifier.
        \item \textbf{Refus}~: la demande est archivée avec le statut
              «~refusée~».
        \item \textbf{Suppression}~: la demande est définitivement effacée
              de la base de données.
        \item \textbf{Champs non conformes}~: aucune demande n'est créée~;
              l'utilisateur corrige et ressoumet.
    \end{itemize}

    \item[Exceptions]
    \textbf{Exception : Identifiant déjà utilisé}
    \begin{itemize}[leftmargin=*, nosep]
        \item Si l'identifiant est déjà attribué à un compte existant ou à
              une demande en attente, le système affiche
              «~Cet identifiant est déjà utilisé~» et bloque la soumission.
    \end{itemize}
\end{description}

\end{document}
